\section{导论：这篇论文在问什么}\label{sec:intro}

\subsection{五分钟版}
Witten 1989 年的论文做了一件事：\textbf{把 Jones 多项式变成三维流形的内蕴不变量}。
路线只有一句：三维 Chern--Simons 理论不含度规，量子化之后沿曲面 $\Sigma$ 切开，
量子 Hilbert 空间 $\cH_\Sigma$ 恰好是二维共形场论的共形块空间；把纽结画成 Wilson 线放进
这个理论，链环的配分函数就是 Jones 多项式（按 $q=e^{2\pi i/(N+k)}$ 取值）。

这句话里每一步都值得停下来问。本讲义就是为这些“停下来“写的：原文把细节压成
“one can show”、“it is easy to see“、“by Segal”，而各本教材补的又不是同一块，
归一化还彼此冲突。所以每一讲的结构固定为：原文说了什么（带页码）$\to$ 原文没说什么
（缺口 G-x）$\to$ 我们补的推导（D-x）$\to$ 与其他来源对照 $\to$ 陷阱与练习。

\subsection{Atiyah 的两个问题}
1988 年，Atiyah 在 Hermann Weyl 讨论会上提了两个问题 \src{351}，本文开篇即引 \cite{atiyah1988}：
\begin{enumerate}
  \item 为 Donaldson 的四维规范理论不变量找一个物理解释；
  \item 为 Jones 多项式找一个内蕴的三维定义。
\end{enumerate}
第一个问题 Witten 在 1988 年的《Topological quantum field theory》\cite{witten1988tqft} 里
用扭曲 $N=2$ 超对称 Yang--Mills 解决了；本文解决第二个。注意原文只说
“第二问是本文的主题“（\emph{the second question is the subject of the present paper}），
并没有说“本文证明了某某定理“——摘要用的动词是 \emph{gives a natural framework}、
\emph{can be generalized}、\emph{computable} \src{351}。这个语气差别是全书的第一条纪律：
\textbf{分清“作者断言“与“作者证明“}。

\subsection{Jones 多项式缺什么}
Jones 1983 年的定义是二维的：取纽结的投影图，在每个交叉处用
$t^{\pm1}$ 做拆接关系（skein relation），再归一到平凡纽结 \cite{jones1985}。
这定义优雅、可算，但\textbf{依赖投影图}：换一个投影，中间步骤全变。
\begin{keypoint}[原文的类比]
Witten 的原话 \src{352}：1930 年代人们知道量子电动力学是相对论性的，
却不知道一个明显相对论性的表述；后来有了量子场论才补上。
Jones 多项式的处境一样：定义是二维的，但“是什么东西使它成为三维不变量“一直不清楚。
\end{keypoint}
补上之后要付的代价是：Jones 多项式变成三维理论的\textbf{配分函数}，
于是它自然推广到任意三维流形上的链环，也自然带上 Wilson 线的表示标签。

\subsection{物理侧的动机：二维来源太多}
纽结多项式在 1980 年代从至少三条互不相干的路线冒出来 \src{352}：
可解格点模型（Kauffman、Turaev）\cite{kauffman1987,turaev1988}、
Yang--Baxter 方程、以及二维共形场论的单值性（Tsuchiya--Kanie）\cite{tuj}。
Witten 的问题意识是：\textbf{为这种多样性找到一个统一主题}，
并弄清“二维共形场论里究竟是什么东西是三维的“。
答案就是 Chern--Simons 理论——而这也是本文对二维 CFT 的“出人意料的启示“。

\subsection{“拓扑不变性 = 广义协变”}
这是全文的立足点 \src{353}。通常说一个理论“拓扑不变“，是指配分函数不依赖度规；
广义相对论里“对度规积分“只是实现这种不变性的一种方式。
Chern--Simons 给了另一种方式：\textbf{作用量本身不含度规}，所以量子化出来自然自动不依赖度规。
\begin{keypoint}[读法]
这句话是“为什么是 Chern--Simons”的全部理由：不是因为 CS 简单，
而是因为它是三维上最一般的、只依赖定向的、规范不变的局域作用量（在低能展开的第一阶）。
凡是含度规的作用量（Yang--Mills、Einstein--Hilbert）量子化后都会得到依赖度规的量。
\end{keypoint}

\subsection{本文的论证结构：四层，不是一层}
\begin{center}
\begin{tabular}{@{}clp{8.2cm}@{}}
\toprule
层 & 原文位置 & 内容 \\
\midrule
L1 构造 & \S1 & 写出作用量 (1.3)，说明 $k\in\mathbb Z$，定义 Wilson 线与配分函数 (1.6) \\
L2 自洽 & \S2--\S3 & 弱耦合展开与正则量子化两条互补路线；给出 $\cH_\Sigma$ 与共形块空间的等同 \\
L3 计算 & \S4 & 手术公式 (4.28)，把任意流形的配分函数化成 $S^3$ 上的链环不变量；(4.51) 给出 Verlinde 公式 \\
L4 解释 & \S5 & 把三维量子化读成二维 CFT 的“相容性条件的来源“ \\
\bottomrule
\end{tabular}
\end{center}
四层的证据强度不同：L1 是定义级（但 $k\in\mathbb Z$ 的“充分性“原文只说必要）；
L2 里“等同“只在对象级被论证，函子性只到 p.397 的 “It is obvious at this point”；
L3 是计算方案 + 部分证明；L4 是解释。\textbf{把 L2 当成“证明了两个理论等价“是本讲义要纠正的最常见误读。}

\subsection{本讲义的读法}
\begin{enumerate}
  \item \textbf{第一遍（一天）}：第 \ref{sec:intro}、\ref{sec:cs}、\ref{sec:verlinde} 讲，
        抓住“作用量 $\to$ 相空间 $\to$ 共形块 $\to$ 融合规则“四步，不追细节。
  \item \textbf{第二遍（两三天）}：第 \ref{sec:phase}--\ref{sec:conformal} 讲与第 \ref{sec:surgery} 讲，
        动笔复算 (3.1)(3.2)(4.28)(4.51)；每道练习都做。
  \item \textbf{第三遍（按需）}：回到具体公式核对约定（$k$ 的归一化、$\mathrm{Tr}$ 的表示、半形式修正），
        并区分“原文说 / 我们补 / 后世“。
\end{enumerate}

\subsection{一张地图：十二讲对应原文哪里}
\begin{center}
\begin{tabular}{@{}lll@{}}
\toprule
讲 & 主题 & 原文位置 \\
\midrule
第 2 讲 & Chern--Simons 作用量与 $k$ 的整性 & \S1，p.353--356 \\
第 3 讲 & 相空间、高斯约束、辛商 & \S3 开头，p.366--368 \\
第 4 讲 & 预量子化、极化、半形式 & \S3.1，p.368--369（原文极简） \\
第 5 讲 & 共形块与射影平坦性 & \S3.2，p.369--371 \\
第 6 讲 & Wilson 线、穿刺点、Borel--Weil--Bott & \S3.3--3.4，p.371--374 \\
第 7 讲 & WZW 与共形场论工具箱 & 原文只引用，未展开 \\
第 8 讲 & 手术、$S/T$ 矩阵、标架反常 & \S4.1--4.5，p.376--393 \\
第 9 讲 & Verlinde 公式：推导与地位 & \S4.5，p.388--393 \\
第 10 讲 & CS/WZW 与 AdS$_3$/CFT$_2$：两种对偶 & \S5，p.394--397 \\
第 11 讲 & 纽结同调的范畴化（延伸） & 原文无，后世 \\
第 12 讲 & 边界、开放问题、三遍读法 & \S5 末，p.397 \\
\bottomrule
\end{tabular}
\end{center}

\subsection{如果你只有一小时}
读第 2 讲的 (1.3)(1.4) 与第 9 讲的十步链，然后做第 9 讲练习 1（$SU(2)_{k=1}$ 的 $S$ 矩阵复算）。
做完你会知道：这篇论文的核心不是“某个深刻定理“，而是一个\textbf{可精确求解的构造}
加上一整套自洽的计算，以及若干断言。把它当成“已经严格的定理集“会处处别扭；
把它当成“物理学家给出的、后来被数学严格化的纲领“，每一节的位置就都清楚了。

\begin{boundary}
\begin{itemize}
  \item 已证（原文或引用）：$k\in\mathbb Z$ 的必要性；(1.6) 在 $S^3$ 上给出 Jones（断言+弱耦合验证）；
        $\cH_\Sigma$ 有限维；手术公式在若干例子上的自洽。
  \item 我们补：$k\in\mathbb Z$ 的充分性论证（第 2 讲 D-2.4）；$\dim\cH$ 的半经典计数（第 4 讲）；
        (4.51) 的十步链（第 9 讲）。
  \item 未证/留白：函子性（p.397）、一般手术下的标架追踪（p.391）、半形式修正（原文未做）、
        量子群构造（p.372 只有一句话）。
\end{itemize}
\end{boundary}
